% Historical source (removed by author request): references/_OLD_Übungsaufgaben/Übungsblatt 4 - Funktionen/Blatt4.tex:234
\begin{ExerciseBlock}
\ExerciseHeading{14}{Trigonometrische Zusammenhänge}{functions_trigonometric_identities}
Rechnen Sie mit den Additionstheoremen
\begin{equation*}
	\begin{split}
		\sin(\alpha \pm \beta) &= \sin(\alpha) \cos(\beta) \pm \sin(\beta) \cos(\alpha) \\
		\cos(\alpha \pm \beta) &= \cos(\alpha) \cos(\beta) \mp \sin(\alpha) \sin(\beta)
	\end{split}
\end{equation*}
die folgenden Beziehungen nach:
\begin{enumerate}[label=\alph*)]
	\item $\displaystyle \sin(2x) = 2 \sin(x) \cos(x)$
	\item $\displaystyle \cos(2x) = 2 \cos^2(x) - 1$
	\item $\displaystyle \sin(x + 2 \pi) = \sin(x)$
	\item $\displaystyle \sin(x \pm \pi) = -\sin(x)$
\end{enumerate}

\end{ExerciseBlock}
\UseExerciseSolution{functions_trigonometric_identities}
